% 防退化说明：公开随笔样张必须保持通用，不得放入作者实际稿件、出版物专属
% 文案或未公开书目；它只验证 essay profile、题名页和叙事照片接口。
\documentclass[profile=essay,titlelayout=essay]{atelier}

\title{从一张照片开始的城市观察}
\subtitle{把场景、材料与论证组织成可以连续阅读的长文}
\author{AtelierTeX Project · 54wsdf}
\kicker{ESSAY / PUBLIC SPECIMEN}
\publicationdate{2026-08-20}
\AtelierSetAffiliation{AtelierTeX 通用排版样张}
\AtelierSetDocumentMeta{ESSAY / ATELIERTEX 0.5}
\AtelierSetAuthorNote{本页为通用功能示例，不对应任何真实投稿稿件。}
\AtelierDeck{随笔排版不是放弃证据结构，而是让题名、阅读入口、叙事照片与正文论证在同一套长文语义中获得更自然的节奏。}

\begin{document}
\maketitle

\section{观察从哪里开始}

一张照片可以先承担场景说明，再成为问题的入口。作者需要区分照片正在呈现什么、观察者据此提出了什么问题，以及哪些判断仍需要文献、访谈或数据支持。排版只负责建立层级，不能代替证据本身。

\AtelierPhotoPlaceholder[34mm]{叙事照片位置}{用于交代地点、时间与观察对象的通用照片说明。}{PHOTO / TO BE SELECTED}

\section{从场景走向论证}

当材料进入正文，章节结构仍然负责推进问题、证据与解释。叙事照片不编号，适合提供现场感；需要在正文中精确引用的分析图、统计图和截图证据继续使用正式编号图件。

这种分工让文化随笔、研究札记和公开长文可以共享同一套多语种、图表与书目基础设施，同时保留各自的阅读节奏。

\end{document}
